\documentclass[aspectratio=169,10pt]{beamer}

\usepackage{iftex}
\providecommand{\dunenoirsetupfonts}{}
\newcommand{\DuneNoirLocalFontPath}{fonts}

\newcommand{\DuneNoirMonoFallback}{%
  \IfFontExistsTF{JetBrainsMono Nerd Font Propo}{
    \setmonofont{JetBrainsMono Nerd Font Propo}[Scale=0.92]
  }{
    \IfFontExistsTF{JetBrainsMono Nerd Font}{
      \setmonofont{JetBrainsMono Nerd Font}[Scale=0.92]
    }{
      \setmonofont{Courier New}[Scale=0.92]
    }
  }
}

\newcommand{\DuneNoirSansFallback}{%
  \IfFontExistsTF{IBM Plex Sans}{
    \setsansfont{IBM Plex Sans}
    \setmainfont{IBM Plex Sans}
  }{
    \IfFontExistsTF{Helvetica Neue}{
      \setsansfont{Helvetica Neue}
      \setmainfont{Helvetica Neue}
    }{
      \setsansfont{Arial}
      \setmainfont{Arial}
    }
  }
}

\newcommand{\DuneNoirUseHelvetica}{%
  \renewcommand{\dunenoirsetupfonts}{%
    \IfFontExistsTF{Helvetica Neue}{
      \setsansfont{Helvetica Neue}
      \setmainfont{Helvetica Neue}
    }{
      \IfFontExistsTF{Helvetica}{
        \setsansfont{Helvetica}
        \setmainfont{Helvetica}
      }{
        \DuneNoirSansFallback
      }
    }
    \DuneNoirMonoFallback
  }%
}

\newcommand{\DuneNoirUseOpenSans}{%
  \renewcommand{\dunenoirsetupfonts}{%
    \IfFileExists{\DuneNoirLocalFontPath/OpenSans.ttf}{
      \setsansfont{Open Sans Custom}[
        Path=\DuneNoirLocalFontPath/,
        UprightFont={OpenSans.ttf},
        ItalicFont={OpenSans-Italic.ttf},
        BoldFont={OpenSans.ttf},
        BoldItalicFont={OpenSans-Italic.ttf},
        BoldFeatures={FakeBold=1.45},
        BoldItalicFeatures={FakeBold=1.45}
      ]
      \setmainfont{Open Sans Custom}[
        Path=\DuneNoirLocalFontPath/,
        UprightFont={OpenSans.ttf},
        ItalicFont={OpenSans-Italic.ttf},
        BoldFont={OpenSans.ttf},
        BoldItalicFont={OpenSans-Italic.ttf},
        BoldFeatures={FakeBold=1.45},
        BoldItalicFeatures={FakeBold=1.45}
      ]
    }{
      \DuneNoirSansFallback
    }
    \DuneNoirMonoFallback
  }%
}

\newcommand{\DuneNoirUseDosis}{%
  \renewcommand{\dunenoirsetupfonts}{%
    \IfFileExists{\DuneNoirLocalFontPath/Dosis.ttf}{
      \setsansfont{Dosis Custom}[
        Path=\DuneNoirLocalFontPath/,
        UprightFont={Dosis.ttf},
        ItalicFont={Dosis.ttf},
        BoldFont={Dosis.ttf},
        BoldItalicFont={Dosis.ttf},
        ItalicFeatures={FakeSlant=0.16},
        BoldFeatures={FakeBold=1.6},
        BoldItalicFeatures={FakeBold=1.6,FakeSlant=0.16}
      ]
      \setmainfont{Dosis Custom}[
        Path=\DuneNoirLocalFontPath/,
        UprightFont={Dosis.ttf},
        ItalicFont={Dosis.ttf},
        BoldFont={Dosis.ttf},
        BoldItalicFont={Dosis.ttf},
        ItalicFeatures={FakeSlant=0.16},
        BoldFeatures={FakeBold=1.6},
        BoldItalicFeatures={FakeBold=1.6,FakeSlant=0.16}
      ]
    }{
      \DuneNoirSansFallback
    }
    \DuneNoirMonoFallback
  }%
}

\newcommand{\DuneNoirUseKoHo}{%
  \renewcommand{\dunenoirsetupfonts}{%
    \IfFileExists{\DuneNoirLocalFontPath/KoHo-Regular.ttf}{
      \setsansfont{KoHo Custom}[
        Path=\DuneNoirLocalFontPath/,
        UprightFont={KoHo-Regular.ttf},
        BoldFont={KoHo-Bold.ttf},
        ItalicFont={KoHo-Italic.ttf},
        BoldItalicFont={KoHo-BoldItalic.ttf}
      ]
      \setmainfont{KoHo Custom}[
        Path=\DuneNoirLocalFontPath/,
        UprightFont={KoHo-Regular.ttf},
        BoldFont={KoHo-Bold.ttf},
        ItalicFont={KoHo-Italic.ttf},
        BoldItalicFont={KoHo-BoldItalic.ttf}
      ]
    }{
      \DuneNoirSansFallback
    }
    \DuneNoirMonoFallback
  }%
}

\DuneNoirUseHelvetica
\makeatletter
\input{../dune_dark_template_lab/beamerthemeDUNENoir.sty}
\makeatother

\setbeamertemplate{headline}{}
\setbeamerfont{frametitle}{size=\Large,series=\bfseries}
\setbeamertemplate{frametitle}{%
  \vspace*{0.8mm}%
  {\usebeamerfont{frametitle}\usebeamercolor[fg]{frametitle}\insertframetitle\par}%
  \vspace{0.7mm}%
  {\color{duneorange}\rule{0.13\paperwidth}{1.1pt}\par}%
  \vspace{0.4mm}%
}
\setbeamertemplate{footline}{%
  \leavevmode%
  \begin{beamercolorbox}[
    wd=\paperwidth,
    ht=4.6mm,
    dp=1.0mm,
    leftskip=0.8cm,
    rightskip=0.8cm
  ]{footline}%
    {\tiny\textcolor{noirmuted}{\talkdate}}%
    \hfill%
    {\tiny\textcolor{duneorange}{\bfseries\insertframenumber}%
      \textcolor{noirmuted}{/\inserttotalframenumber}}%
  \end{beamercolorbox}%
}

\ifPDFTeX
  \usepackage[T1]{fontenc}
  \usepackage[utf8]{inputenc}
\fi

\usepackage{amsmath,amssymb}
\usepackage{array,booktabs,tabularx}
\usepackage{graphicx}
\usepackage{hyperref}
\usepackage{tikz}
\usetikzlibrary{
  arrows.meta,
  calc,
  fit,
  positioning,
  shapes.geometric,
  backgrounds
}

\title[Executable SC/DPS Plan]{From Interfaces to an\\Executable SC/DPS Plan}
\subtitle{A simulation and verification framework\\that turns system interactions\\into the development schedule}
\author{Manuel Arroyave}
\date{SC/DPS coordination discussion --- 29 July 2026}
\newcommand{\mydate}{29 July 2026}

\definecolor{evidencegreen}{HTML}{8FD6A9}
\definecolor{projectioncyan}{HTML}{75C9D6}
\definecolor{questionorange}{HTML}{F4A261}
\definecolor{riskred}{HTML}{F28B82}
\definecolor{warmyellow}{HTML}{E7C547}
\definecolor{purpletest}{HTML}{B7A4E8}

\setbeamercolor{panel evidence}{fg=noirbg,bg=evidencegreen}
\setbeamercolor{panel projection}{fg=noirbg,bg=projectioncyan}
\setbeamercolor{panel unresolved}{fg=noirbg,bg=questionorange}
\setbeamercolor{panel decision}{fg=noirbg,bg=warmyellow}

\setbeamertemplate{title page}{%
  \begin{tikzpicture}[remember picture,overlay]
    \fill[noirbg] (current page.south west) rectangle (current page.north east);
    \fill[noirpanelalt] (current page.north west)
      rectangle ($(current page.north east)+(0,-3.15)$);
    \fill[duneorange] ($(current page.north west)+(0.82,-0.72)$)
      rectangle ++(1.55,0.08);
    \fill[noircyan] ($(current page.south east)+(-2.8,0.8)$)
      rectangle ++(1.9,0.08);
    \IfFileExists{../dune_dark_template_lab/logos/DUNElogo_white.png}{%
      \node[anchor=south east,opacity=0.13]
        at ($(current page.south east)+(-0.55,1.45)$) {%
        \includegraphics[width=0.33\paperwidth]
          {../dune_dark_template_lab/logos/DUNElogo_white.png}%
      };
    }{}%
    \IfFileExists{../dune_dark_template_lab/logos/FNAL-Logo-NAL-Light.png}{%
      \node[anchor=south east]
        at ($(current page.south east)+(-0.92,0.82)$) {%
        \includegraphics[height=0.45cm]
          {../dune_dark_template_lab/logos/FNAL-Logo-NAL-Light.png}%
      };
    }{}%
    \node[
      anchor=north west,
      text width=0.78\paperwidth,
      text=white,
      font=\bfseries\fontsize{25}{28}\selectfont
    ] at ($(current page.north west)+(0.82,-0.96)$) {\inserttitle};
    \node[
      anchor=north west,
      text width=0.54\paperwidth,
      text=noirmuted,
      font=\fontsize{13}{17}\selectfont
    ] at ($(current page.north west)+(0.82,-3.78)$) {\insertsubtitle};
    \node[
      anchor=north west,
      text=noircyan,
      font=\bfseries
    ] at ($(current page.north west)+(0.82,-6.55)$) {\insertauthor};
    \node[
      anchor=north west,
      text=noirmuted,
      font=\small
    ] at ($(current page.north west)+(0.82,-7.08)$) {\insertdate};
  \end{tikzpicture}
  \vfill
}

\newcommand{\eyebrow}[1]{{\color{noircyan}\scriptsize\bfseries #1\par}}
\newcommand{\tinycite}[1]{{\color{noirmuted}\tiny #1\par}}
\newcommand{\tightitem}{%
  \setlength{\itemsep}{0.8mm}%
  \setlength{\parskip}{0pt}%
}

\newcommand{\panelbox}[3][panel]{%
  \begin{beamercolorbox}[wd=\linewidth,sep=5pt,rounded=true]{#1}
    {\usebeamercolor[fg]{#1}\bfseries #2\par}
    \vspace{0.55mm}
    {\usebeamercolor[fg]{#1}\footnotesize #3\par}
  \end{beamercolorbox}
}

\newcommand{\fixedpanel}[4][panel]{%
  \begin{beamercolorbox}[wd=\linewidth,sep=4pt,rounded=true]{#1}
    \begin{minipage}[t][#2][t]{\dimexpr\linewidth-8pt\relax}
      \raggedright
      {\usebeamercolor[fg]{#1}\bfseries #3\par}
      \vspace{0.45mm}
      {\usebeamercolor[fg]{#1}\footnotesize #4\par}
    \end{minipage}
  \end{beamercolorbox}
}

\newcommand{\evidencepill}[2]{%
  \tikz[baseline=-0.55ex]\node[
    rounded corners=3pt,
    draw=#1,
    fill=#1!18!noirbg,
    inner xsep=4pt,
    inner ysep=1.6pt,
    text=#1,
    font=\tiny\bfseries
  ]{#2};%
}

\newcommand{\statcard}[3][panel alt]{%
  \begin{beamercolorbox}[wd=\linewidth,sep=3pt,rounded=true]{#1}
    {\color{noircyan}\bfseries\fontsize{12}{14}\selectfont #2\par}
    \vspace{0.2mm}
    {\color{noirmuted}\tiny #3\par}
  \end{beamercolorbox}
}

\newcommand{\sectioncard}[2]{%
  \begin{frame}[plain]
    \begin{tikzpicture}[remember picture,overlay]
      \fill[noirpanelalt] (current page.south west)
        rectangle (current page.north east);
      \fill[duneorange] ($(current page.north west)+(0.9,-0.9)$)
        rectangle ++(1.8,0.08);
      \fill[noircyan] ($(current page.south east)+(-2.8,0.9)$)
        rectangle ++(1.8,0.08);
    \end{tikzpicture}
    \vspace*{0.21\paperheight}
    {\color{noirmuted}\small\bfseries #1\par}
    \vspace{3mm}
    {\color{white}\bfseries\fontsize{28}{31}\selectfont #2\par}
  \end{frame}
}

\tikzset{
  systembox/.style={
    draw=noirline,
    fill=noirpanelalt,
    text=noirtext,
    rounded corners=4pt,
    align=center,
    text width=2.25cm,
    minimum height=0.9cm,
    inner sep=4pt,
    font=\scriptsize
  },
  systemcyan/.style={
    systembox,
    draw=noircyan,
    fill=noircyan!16!noirpanel
  },
  systemmint/.style={
    systembox,
    draw=evidencegreen,
    fill=evidencegreen!15!noirpanel
  },
  systemorange/.style={
    systembox,
    draw=questionorange,
    fill=questionorange!15!noirpanel
  },
  systempurple/.style={
    systembox,
    draw=purpletest,
    fill=purpletest!14!noirpanel
  },
  flowarrow/.style={-Latex,draw=noircyan,line width=0.85pt},
  interfaceline/.style={draw=noircyan,line width=0.85pt},
  evidencearrow/.style={
    -Latex,
    draw=purpletest,
    dashed,
    line width=0.85pt
  },
  evidenceline/.style={
    draw=purpletest,
    dashed,
    line width=0.85pt
  },
  warnarrow/.style={-Latex,draw=questionorange,line width=0.9pt},
  riskarrow/.style={-Latex,draw=riskred,line width=1.0pt},
  labelnote/.style={
    font=\tiny,
    text=noirmuted,
    align=center,
    fill=noirbg,
    inner sep=2pt
  },
  stagebox/.style={
    draw=noirline,
    fill=noirpanelalt,
    text=noirtext,
    rounded corners=4pt,
    align=left,
    text width=2.55cm,
    minimum height=1.55cm,
    inner sep=5pt,
    font=\scriptsize
  },
  decisionbox/.style={
    draw=warmyellow,
    fill=warmyellow!14!noirpanel,
    text=noirtext,
    rounded corners=5pt,
    align=left,
    text width=2.5cm,
    minimum height=1.65cm,
    inner sep=6pt,
    font=\scriptsize
  }
}

\begin{document}

{
\setbeamertemplate{headline}{}
\setbeamertemplate{footline}{}
\begin{frame}[plain,noframenumbering]
  \titlepage
\end{frame}
}

\begin{frame}[t]{This proposal connects existing work}
\eyebrow{INTEGRATION, NOT A RESET}
\vspace{1.2mm}

\begin{columns}[T,onlytextwidth]
\column{0.48\textwidth}
\fixedpanel{4.0cm}{What the collaboration already has}{
\begin{itemize}\tightitem
  \item Years of subsystem design, firmware, controls, and test-stand work.
  \item Working protocols, serializers, device servers, ICDs, and DAQ packages.
  \item Experienced consortium and integration teams who know their hardware.
  \item Existing SC and DPS decisions that must remain traceable.
\end{itemize}
}

\column{0.48\textwidth}
\fixedpanel{4.0cm}{The missing shared artifact}{
\begin{itemize}\tightitem
  \item One lifecycle-aware view of interactions across subsystem boundaries.
  \item Explicit ownership of normal control, protection, failure, and recovery.
  \item An executable link from interface requirements to package-level tests.
  \item A development schedule derived from real dependencies and risks.
\end{itemize}
}
\end{columns}

\vspace{1.2mm}
\begin{beamercolorbox}[wd=\linewidth,sep=6pt,rounded=true]{panel alt}
\centering
\color{white}\bfseries
Reuse the subsystem work. Connect it into one versioned, testable system contract.
\end{beamercolorbox}
\end{frame}

\begin{frame}[t]{Why move now?}
\eyebrow{INTERFACE ASSUMPTIONS ARE BECOMING PHYSICAL CONSTRAINTS}
\vspace{2mm}

\begin{tikzpicture}[x=1cm,y=1cm]
  \node[stagebox,draw=evidencegreen] (design) at (1.45,0) {
    \textcolor{evidencegreen}{\bfseries 1\enspace Design}\\[1mm]
    Protocols, power architecture, authority, states, timing.
  };
  \node[stagebox] (procure) at (4.35,0) {
    \textcolor{noircyan}{\bfseries 2\enspace Procurement}\\[1mm]
    PLCs, PDUs, servers, switches, UPS, field interfaces.
  };
  \node[stagebox,draw=questionorange] (install) at (7.25,0) {
    \textcolor{questionorange}{\bfseries 3\enspace Installation}\\[1mm]
    Cabling, rack maps, network zones, interlock paths.
  };
  \node[stagebox,draw=riskred] (commission) at (10.15,0) {
    \textcolor{riskred}{\bfseries 4\enspace Commissioning}\\[1mm]
    Recovery, partial states, fault response, acceptance evidence.
  };
  \node[stagebox,draw=riskred] (operate) at (12.93,0) {
    \textcolor{riskred}{\bfseries 5\enspace Operation}\\[1mm]
    Availability, change control, incident response.
  };

  \draw[line width=6pt,evidencegreen!50]
    (0.25,-1.42) -- (4.85,-1.42);
  \draw[line width=6pt,questionorange!65]
    (4.85,-1.42) -- (9.45,-1.42);
  \draw[line width=6pt,riskred!75]
    (9.45,-1.42) -- (13.95,-1.42);
  \node[font=\tiny,text=noirbg] at (2.55,-1.42) {cheap to change};
  \node[font=\tiny,text=noirbg] at (7.15,-1.42) {coordination cost grows};
  \node[font=\tiny,text=noirbg] at (11.75,-1.42) {hardware/recovery rework};

  \draw[draw=noircyan,line width=1.0pt] (0.25,-1.88) -- (13.52,-1.88);
  \fill[noircyan]
    (13.95,-1.88) -- (13.52,-1.64) -- (13.52,-2.12) -- cycle;
  \node[labelnote,above] at (7.1,-1.88)
    {project progression and increasing cost of change};
\end{tikzpicture}

\vspace{3mm}
\begin{columns}[T,onlytextwidth]
\column{0.48\textwidth}
\panelbox{Move fast on unknowns}{
Resolve ownership, state constraints, sequencing, and independent protection
before they are frozen into hardware and deployment choices.}
\column{0.48\textwidth}
\panelbox{Do not rush protection}{
Speed means earlier questions and executable evidence---not bypassing safety,
consortium, electrical, or configuration reviews.}
\end{columns}
\end{frame}

\begin{frame}[t]{A draft bench already exists}
\eyebrow{SEPARATE NATIVE PATHS; ONE OPERATOR AND TEST VIEW}
\vspace{1mm}

\begin{center}
\resizebox{0.88\linewidth}{!}{%
\begin{tikzpicture}[x=1cm,y=1cm]
  \node[systemcyan,text width=2.75cm] (boards) at (0,1.45) {
    190 DAPHNE + 480 WIB/PTC\\board services
  };
  \node[systemmint,text width=2.75cm] (pdu) at (0,0.48) {
    108 PDUs\\SNMP endpoints
  };
  \node[systemmint,text width=2.75cm] (power) at (0,-0.49) {
    208 power channels\\SCPI / native OPC-UA
  };
  \node[systemorange,text width=2.75cm] (env) at (0,-1.46) {
    sensors, permits,\\alarms
  };

  \node[systemcyan,text width=3.20cm] (fleet) at (4.05,1.45) {
    DAPHNE/WIB OPC-UA fleet\\shared bridge implementation
  };
  \node[systemmint,text width=3.20cm] (snmp) at (4.05,0.48) {
    SNMP driver / bridge
  };
  \node[systemmint,text width=3.20cm] (supply) at (4.05,-0.49) {
    supply driver / adapter
  };
  \node[systemorange,text width=3.20cm] (plc) at (4.05,-1.46) {
    PLC/DPS or native OPC-UA\\hard action remains local
  };

  \node[systemorange,text width=3.25cm] (daq) at (8.65,2.55) {
    DUNE DAQ native clients\\
    \texttt{drunc}/appfwk + opmon
  };
  \node[systemmint,text width=3.25cm,minimum height=1.05cm] (ignition) at (8.65,0) {
    Ignition Gateway on Linux\\
    tags, alarms, history, navigation
  };

  \draw[interfaceline] (boards) -- (fleet);
  \draw[interfaceline] (pdu) -- (snmp);
  \draw[interfaceline] (power) -- (supply);
  \draw[interfaceline] (env) -- (plc);
  \draw[interfaceline] (boards.north east) to[out=24,in=180]
    node[labelnote,above]{native ZMQ/protobuf} (daq.west);

  \draw[interfaceline] (6.35,-1.46) -- (6.35,1.45);
  \foreach \adapter in {fleet,snmp,supply,plc}
    \draw[interfaceline] (\adapter.east) -- (6.35,0 |- \adapter.east);
  \draw[interfaceline] (6.35,0) -- (ignition.west);

  \node[labelnote,text width=2.9cm] at (6.35,1.88)
    {SC provider boundary};
  \node[labelnote,text width=3.3cm] at (8.65,1.38)
    {unheaded lines = interfaces};
  \node[labelnote,text=purpletest,text width=3.6cm] at (8.65,-1.18)
    {evidence harness: timing, faults, readback, provenance};
\end{tikzpicture}
}
\end{center}

\vspace{0.2mm}
\tinycite{\centering
Current synthetic inventory: 1,818 logical objects \textbullet\
40,755 published variables \textbullet\
986 writable endpoints \textbullet\
one shared Ignition operator layer.}

\end{frame}

\begin{frame}[t]{Evidence must remain labelled}
\eyebrow{WHAT WAS MEASURED, WHAT WAS PROJECTED, WHAT IS STILL OPEN}
\vspace{1.6mm}

\begin{columns}[T,onlytextwidth]
\column{0.32\textwidth}
\fixedpanel[panel evidence]{4.93cm}{Tested locally}{
\begin{itemize}\tightitem
  \item 12,000 balanced configure/readback trials; zero observed errors.
  \item One 40,755-variable OPC-UA snapshot in 0.516 s.
  \item One local pass over 986 writable endpoints in 0.517 s.
  \item Four same-host OPC-UA processes at the measured monitoring knee.
\end{itemize}
\vfill
\evidencepill{evidencegreen}{T = workstation evidence}
}

\column{0.32\textwidth}
\fixedpanel[panel projection]{4.93cm}{Source-derived / projected}{
\begin{itemize}\tightitem
  \item DAPHNE explicit-wait floor: 2.6901 s per board.
  \item Four-FEMB WIB explicit-wait floor: 16.01--16.10 s.
  \item \texttt{drunc}/appfwk fan-out and queue models.
  \item Bridge/leaf sizing with posterior timeout and retry uncertainty.
\end{itemize}
\vfill
\evidencepill{projectioncyan}{P = planning model}
}

\column{0.32\textwidth}
\fixedpanel[panel unresolved]{4.93cm}{Not yet demonstrated}{
\begin{itemize}\tightitem
  \item Real PLCs, switches, detector networks, and production endpoints.
  \item Hardware I/O and calibration time beyond explicit source waits.
  \item Full Ignition subscription, historian, alarm, and client load.
  \item Safety integrity, physical interlocks, failover, or production sizing.
\end{itemize}
\vfill
\evidencepill{questionorange}{U = owner input / test needed}
}
\end{columns}

\vspace{0mm}
\centering
\color{noirmuted}\scriptsize
The bench is useful because it separates these categories instead of turning a local result into a detector guarantee.
\end{frame}

\begin{frame}[t]{Monitoring is not the sizing workload}
\eyebrow{CONFIGURATION, QUEUES, AND RECOVERY SET THE CRITICAL PATH}
\vspace{1.0mm}

\begin{columns}[T,onlytextwidth]
\column{0.61\textwidth}
\begin{tikzpicture}[x=0.125cm,y=0.76cm]
  \node[anchor=east,font=\scriptsize,text=noirmuted] at (0,4.1) {OPC-UA snapshot};
  \draw[fill=noircyan!65,draw=noircyan,rounded corners=1pt]
    (0.3,3.85) rectangle (0.82,4.35);
  \node[anchor=west,font=\scriptsize,text=noirtext] at (1.05,4.1)
    {0.516 s \quad 40,755 variables, local software path};

  \node[anchor=east,font=\scriptsize,text=noirmuted] at (0,3.1) {DAPHNE board};
  \draw[fill=evidencegreen!55,draw=evidencegreen,rounded corners=1pt]
    (0.3,2.85) rectangle (3.0,3.35);
  \node[anchor=west,font=\scriptsize,text=noirtext] at (3.22,3.1)
    {2.69 s \quad explicit source waits};

  \node[anchor=east,font=\scriptsize,text=noirmuted] at (0,2.1)
    {190-DAPHNE fleet};
  \draw[fill=evidencegreen!72,draw=evidencegreen,rounded corners=1pt]
    (0.3,1.85) rectangle (3.0,2.35);
  \node[anchor=west,font=\scriptsize,text=noirtext] at (3.22,2.1)
    {2.69 s \quad one projected wave};

  \node[anchor=east,font=\scriptsize,text=noirmuted] at (0,1.1) {WIB board};
  \draw[fill=questionorange!60,draw=questionorange,rounded corners=1pt]
    (0.3,0.85) rectangle (16.4,1.35);
  \node[anchor=west,font=\scriptsize,text=noirtext] at (16.62,1.1)
    {16.10 s};

  \node[anchor=east,font=\scriptsize,text=noirmuted] at (0,0.1) {480-WIB fleet};
  \draw[fill=riskred!62,draw=riskred,rounded corners=1pt]
    (0.3,-0.15) rectangle (48.6,0.35);
  \node[anchor=east,font=\scriptsize\bfseries,text=noirbg] at (48.1,0.1)
    {48.30 s};

  \draw[draw=noirline] (0.3,-0.55) -- (50.2,-0.55);
  \foreach \x/\label in {0/0,10/10,20/20,30/30,40/40,50/50} {
    \draw[draw=noirline] ({0.3+\x},-0.48) -- ({0.3+\x},-0.62);
    \node[font=\tiny,text=noirmuted] at ({0.3+\x},-0.82) {\label s};
  }
\end{tikzpicture}

\vspace{1mm}
\tinycite{Fleet rows are no-failure floors for four shared bridge queues with
48 active leaves each: 47--48 DAPHNEs/bridge form one 2.69 s wave; 120
WIBs/bridge form three 16.10 s waves. Both are projections, not hardware
measurements.}

\column{0.36\textwidth}
\fixedpanel{5.2cm}{Why configuration is heavier}{
\begin{itemize}\tightitem
  \item Explicit firmware calibration and verification waits.
  \item DAQ controller/application nesting and bounded fan-out.
  \item Per-bridge queues, timeout, retry, and independent readback.
  \item State, power, and timing prerequisites before a command is legal.
\end{itemize}
\vspace{2mm}
\textcolor{noircyan}{\bfseries Planning point}\\
1 Ignition Gateway\\
4 DAPHNE/WIB bridge instances\\
47--48 DAPHNE + 120 WIB/bridge\\
48 active leaves/bridge
}
\end{columns}
\end{frame}

\begin{frame}[t]{The proposal in one line}
\eyebrow{THE MATRIX PRODUCES THE SCHEDULE}
\vspace{2mm}

\begin{center}
\resizebox{0.97\linewidth}{!}{%
\begin{tikzpicture}[node distance=0.52cm]
  \node[systemcyan,text width=2.35cm,minimum height=1.35cm] (inputs) {
    Controlled\\interface inputs\\[1mm]
    \tiny variables, commands, schemas
  };
  \node[systemorange,text width=2.35cm,minimum height=1.35cm,right=of inputs] (matrix) {
    Stage-aware\\interaction matrix\\[1mm]
    \tiny states, interlocks, recovery
  };
  \node[systemorange,text width=2.35cm,minimum height=1.35cm,right=of matrix] (graph) {
    Dependency +\\hazard graph\\[1mm]
    \tiny critical paths, unknowns
  };
  \node[systemmint,text width=2.35cm,minimum height=1.35cm,right=of graph] (packages) {
    Work packages +\\acceptance tests\\[1mm]
    \tiny owners, evidence gates
  };
  \node[systemmint,text width=2.35cm,minimum height=1.35cm,right=of packages] (schedule) {
    Staffed technical\\schedule\\[1mm]
    \tiny effort, people, dependencies
  };
  \node[systempurple,text width=2.35cm,minimum height=1.35cm,right=of schedule] (evidence) {
    Continuous\\verification evidence\\[1mm]
    \tiny CI, VST, releases
  };
  \draw[flowarrow] (inputs) -- (matrix);
  \draw[warnarrow] (matrix) -- (graph);
  \draw[warnarrow] (graph) -- (packages);
  \draw[flowarrow] (packages) -- (schedule);
  \draw[evidencearrow] (schedule) -- (evidence);
  \coordinate (feedbackright) at ($(evidence.south)+(0,-0.52)$);
  \coordinate (feedbackleft) at ($(matrix.south)+(0,-0.52)$);
  \draw[evidencearrow,rounded corners=3pt]
    (evidence.south) -- (feedbackright) --
    node[labelnote,below] {results update the controlled baseline}
    (feedbackleft) -- (matrix.south);
\end{tikzpicture}
}
\end{center}

\vspace{1.0mm}
\begin{beamercolorbox}[wd=\linewidth,sep=7pt,rounded=true]{panel decision}
\centering
\color{noirbg}\bfseries\fontsize{12}{14}\selectfont
The development schedule is an output of the interaction matrix---\\
not a calendar assembled from independent subsystem task lists.
\end{beamercolorbox}

\vspace{1mm}
\centering
\color{noirmuted}\scriptsize
The matrix sets priority and order. Effort estimates, staffing, procurement, and review capacity set dates.
\end{frame}

\begin{frame}[t]{What must be captured in the matrix?}
\eyebrow{ENOUGH INFORMATION TO TURN EACH INTERACTION INTO WORK AND EVIDENCE}
\vspace{1.5mm}

\begin{columns}[T,onlytextwidth]
\column{0.32\textwidth}
\fixedpanel{3.7cm}{Identity and authority}{
\begin{itemize}\tightitem
  \item Source and destination.
  \item Variable or command; type, unit, range, rate.
  \item Command owner and permitted actor.
  \item Applicable lifecycle stage.
  \item Source ICD/schema and controlled revision.
\end{itemize}
}

\column{0.32\textwidth}
\fixedpanel{3.7cm}{Constraints and response}{
\begin{itemize}\tightitem
  \item Preconditions and allowed detector states.
  \item Permits, inhibits, and interlocks.
  \item Sequence, deadline, timeout, stale limit.
  \item Expected response and independent readback.
  \item Failure, fallback, retry, and safe response.
\end{itemize}
}

\column{0.32\textwidth}
\fixedpanel{3.7cm}{Closure and evidence}{
\begin{itemize}\tightitem
  \item Recovery, reset, rollback, and latching.
  \item Responsible subsystem and integration owners.
  \item Acceptance test and quantitative criterion.
  \item Simulator, VST, hardware, or proof evidence.
  \item Explicit unknown, conflict, and due date.
\end{itemize}
}
\end{columns}

\vspace{1mm}
\begin{beamercolorbox}[wd=\linewidth,sep=6pt,rounded=true]{panel alt}
\centering\small
Every unresolved cell is visible schedule risk; every closed row can become a reviewed work package and regression test.
\end{beamercolorbox}
\end{frame}

\begin{frame}[t]{Stage and state both matter}
\eyebrow{THE CONSTRAINTS ARE NOT ONLY A FUNCTION OF CRYOGENICS}
\vspace{1.3mm}

\begin{columns}[T,onlytextwidth]
\column{0.39\textwidth}
\begin{tikzpicture}[node distance=0.33cm]
  \node[systemcyan,text width=4.75cm] (install) {
    \textbf{Installation}\\partial inventory, temporary services
  };
  \node[systemcyan,text width=4.75cm,below=of install] (integration) {
    \textbf{Integration}\\mixed firmware, incomplete dependencies
  };
  \node[systemorange,text width=4.75cm,below=of integration] (commission) {
    \textbf{Commissioning}\\energization, fault tests, recovery
  };
  \node[systemmint,text width=4.75cm,below=of commission] (operation) {
    \textbf{Operation}\\stable authority, alarms, controlled change
  };
  \node[systemorange,text width=4.75cm,below=of operation] (maintenance) {
    \textbf{Maintenance / recovery}\\bypasses, rollback, requalification
  };
  \draw[flowarrow] (install) -- (integration);
  \draw[flowarrow] (integration) -- (commission);
  \draw[flowarrow] (commission) -- (operation);
  \draw[warnarrow] (operation) -- (maintenance);
\end{tikzpicture}

\column{0.57\textwidth}
\fixedpanel{2.0cm}{State axes that change an allowed action}{
\begin{tabularx}{\linewidth}{@{}>{\bfseries}lX@{}}
Cryogenics & warm, cool-down, stable, abnormal\\
Power & absent, partial, ramping, stable, trip\\
DAQ/timing & boot, initialised, configured, running, fault\\
Network & healthy, degraded, partitioned, stale\\
Authority & local, commissioning, SC, DAQ, DPS, maintenance\\
\end{tabularx}
}

\vspace{1.0mm}
\begin{beamercolorbox}[wd=\linewidth,sep=4pt,rounded=true]{panel alt}
\centering
\color{white}\bfseries\small
Allowed action depends on\\[-0.2mm]
\textcolor{noircyan}{lifecycle \(\cdot\) state \(\cdot\) permits
\(\cdot\) communications \(\cdot\) authority}
\end{beamercolorbox}

\vspace{0.5mm}
\scriptsize
Each consortium must define both normal SC behavior and its DPS/failure
contract at these boundaries. Nominal variables alone do not define the
integration behavior.
\end{columns}
\end{frame}

\begin{frame}[t]{One actor, two clients, independent protection}
\eyebrow{THE LOGICAL CONTRACT BEHIND DAPHNE AND WIB INTEGRATION}
\vspace{1.2mm}

\begin{columns}[T,onlytextwidth]
\column{0.66\textwidth}
\begin{center}
\resizebox{0.99\linewidth}{!}{%
\begin{tikzpicture}[x=1cm,y=1cm]
  \node[systemcyan,text width=2.5cm] (hw) at (0,1.9) {
    Firmware + hardware\\registers, sensors, faults
  };
  \node[systempurple,text width=2.5cm] (actor) at (3.55,1.9) {
    One target-side\\board actor
  };
  \node[systemcyan,text width=2.35cm] (proto) at (7.0,1.9) {
    Canonical\\ZMQ/protobuf
  };
  \node[systemorange,text width=2.55cm] (daq) at (10.5,3.0) {
    DUNE DAQ\\appfwk controller + opmon
  };
  \node[systemmint,text width=2.55cm] (sc) at (10.5,0.8) {
    OPC-UA bridge\\Ignition + history
  };
  \draw[interfaceline] (hw) -- (actor);
  \draw[interfaceline] (actor) -- (proto);
  \draw[interfaceline] (proto) -- (daq);
  \draw[interfaceline] (proto) -- (sc);

  \node[systemorange,text width=2.5cm] (hazard) at (0,-1.2) {
    Hazard sensors\\permits / contacts
  };
  \node[
    systembox,
    draw=riskred,
    fill=riskred!13!noirpanel,
    text width=2.5cm
  ] (dps) at (3.55,-1.2) {
    DPS PLC / local\\protection logic
  };
  \node[
    systembox,
    draw=riskred,
    fill=riskred!13!noirpanel,
    text width=2.35cm
  ] (inhibit) at (7.0,-1.2) {
    Hard inhibit /\\power action
  };
  \draw[riskarrow] (hazard) -- (dps);
  \draw[riskarrow] (dps) -- (inhibit);
  \draw[evidenceline] (dps) -- node[labelnote,below,sloped]{status} (sc);
  \draw[riskarrow,dashed,bend left=20] (inhibit.north) to
    node[labelnote,right]{overtake} (actor.south);
\end{tikzpicture}
}
\end{center}

\column{0.31\textwidth}
\fixedpanel{5.15cm}{Non-negotiable rules}{
\begin{itemize}\tightitem
  \item One process owns AXI, I2C, SPI, DAC, reset, and power transactions.
  \item DAQ and SC may read the same cached status in parallel.
  \item Every write is serialized, state/authority gated, audited, and read back.
  \item \texttt{opmonlib} is a DAQ projection, not a second device schema.
  \item Required protection must survive the failures for which it is credited.
\end{itemize}
}
\end{columns}
\end{frame}

\begin{frame}[t]{Example: power loss exposes a system requirement}
\eyebrow{A CONDITIONAL QUESTION, NOT A CLAIM THAT THE CURRENT DESIGN IS UNSAFE}
\vspace{1mm}

\begin{center}
\resizebox{0.91\linewidth}{!}{%
\begin{tikzpicture}[node distance=0.75cm and 0.85cm]
  \node[systemorange,text width=3.05cm,minimum height=1.0cm] (loss) {
    Facility / rack\\power loss
  };
  \node[systemorange,text width=3.2cm,minimum height=1.0cm,right=1.35cm of loss,yshift=1.05cm] (wib) {
    WIEC/WIB low voltage\\decays
  };
  \node[systemorange,text width=3.2cm,minimum height=1.0cm,right=1.35cm of loss,yshift=-1.05cm] (bias) {
    APA bias system\\shutdown/discharge
  };
  \node[systempurple,text width=3.25cm,minimum height=1.0cm,right=1.2cm of wib,yshift=-1.05cm] (sc) {
    SC/Ignition request\\may disappear with power/network
  };
  \node[systemmint,text width=3.35cm,minimum height=1.0cm,right=1.2cm of sc] (independent) {
    Independent DPS /\\hard-interlock path
  };

  \draw[warnarrow] (loss) -- (wib);
  \draw[warnarrow] (loss) -- (bias);
  \draw[evidencearrow] (loss) -- (sc);
  \draw[draw=riskred,dashed,line width=1pt,<->] (wib) --
    node[labelnote,right,text width=2.0cm]
      {required ordering?\\hold-up times?}
    (bias);
  \draw[riskarrow] (independent) -- (wib);
  \draw[riskarrow] (independent) -- (bias);
  \draw[evidenceline] (sc) -- (independent);

  \node[
    draw=riskred,
    fill=riskred!12!noirpanel,
    text=noirtext,
    rounded corners=4pt,
    align=left,
    text width=13.2cm,
    inner sep=6pt,
    below=1.05cm of sc
  ] {
    \textbf{Question to close:} if WIB power is not UPS-backed, what mechanism
    guarantees the required ordering between APA-bias removal and WIB power loss?
    If software cannot survive the event, the guarantee must be owned elsewhere.
  };
\end{tikzpicture}
}
\end{center}

\vspace{0mm}
\begin{columns}[T,onlytextwidth]
\column{0.32\textwidth}
\panelbox{Inputs needed}{UPS/hold-up, discharge time, field wiring, PLC/DPS logic, network and power failure modes.}
\column{0.32\textwidth}
\panelbox{Matrix output}{Requirement, owner, independent path, allowed states, timing bound, recovery rule.}
\column{0.32\textwidth}
\panelbox{Schedule output}{Design task, integration test, proof test, acceptance gate, unresolved deadline.}
\end{columns}
\end{frame}

\begin{frame}[t]{What must be managed---and where}
\eyebrow{THE CONFIGURATION BASELINE IS LARGER THAN SOURCE CODE}
\vspace{1.2mm}

\begin{columns}[T,onlytextwidth]
\column{0.34\textwidth}
\fixedpanel{4.65cm}{Controlled in Git}{
\begin{itemize}\tightitem
  \item Matrix, interface register, ADRs, risks, ownership.
  \item Canonical schemas, descriptors, compatibility policy.
  \item Firmware/server/DAQ/bridge source and tests.
  \item PLC logic and approved deployment exports.
  \item Deployment templates, release BOM, rollback.
  \item Test definitions and evidence metadata.
\end{itemize}
}

\column{0.30\textwidth}
\fixedpanel{4.65cm}{External sources of record}{
\begin{itemize}\tightitem
  \item EDMS: approved requirements and ICDs.
  \item HWDB: installed identity, topology, connections.
  \item Artifact registry: firmware, packages, containers, checksums.
  \item Secret store: keys, tokens, credentials.
  \item Historian/operations DB: runtime samples and alarm state.
\end{itemize}
}

\column{0.32\textwidth}
\fixedpanel{4.65cm}{Managed as one release}{
\begin{itemize}\tightitem
  \item Firmware/register ABI.
  \item Board service + schema.
  \item DUNE DAQ package + opmon mapping.
  \item OPC-UA bridge + Ignition export.
  \item DPS/permit revision.
  \item Configuration/HWDB export.
  \item Cross-system acceptance evidence.
\end{itemize}
\vspace{1mm}
\textcolor{questionorange}{\bfseries Rule}\\
One compatible bill of materials pins the exact tested set.
}
\end{columns}
\end{frame}

\begin{frame}[t]{Make the repositories enforce the contract}
\eyebrow{COLLABORATION-OWNED GOVERNANCE + EXECUTABLE COMPATIBILITY}
\vspace{1.3mm}

\begin{columns}[T,onlytextwidth]
\column{0.62\textwidth}
\begin{center}
\resizebox{0.98\linewidth}{!}{%
\begin{tikzpicture}[x=1cm,y=1cm]
  \node[
    systemcyan,
    text width=9.45cm,
    minimum height=0.85cm
  ] (group) at (0,3.25) {Collaboration-controlled Git group / namespace};

  \node[systembox,text width=2.75cm] (schema) at (-3.25,1.78) {
    Interface baselines\\ICDs, schemas, matrices
  };
  \node[systembox,text width=2.75cm] (sim) at (0,1.78) {
    Simulators + adapters\\faults, protocol oracles
  };
  \node[systembox,text width=2.75cm] (prod) at (3.25,1.78) {
    Development packages\\firmware, bridges, PLC/SC
  };

  \node[
    systempurple,
    text width=9.45cm,
    minimum height=1.05cm,
  ] (ci) at (0,0.15) {
    Protected CI evidence gate\\
    schema compatibility, contract tests, negative/fault cases, readback,
    version matrix
  };
  \node[
    systemmint,
    text width=9.45cm,
    minimum height=0.85cm,
  ] (release) at (0,-1.25) {
    Reviewed release + signed evidence package
  };

  \draw[interfaceline] (group.south) -- (0,2.58);
  \draw[interfaceline] (-3.25,2.58) -- (3.25,2.58);
  \foreach \repo/\xpos in {schema/-3.25,sim/0,prod/3.25}
    \draw[interfaceline] (\xpos,2.58) -- (\repo.north);

  \draw[evidenceline] (-3.25,1.04) -- (3.25,1.04);
  \foreach \repo/\xpos in {schema/-3.25,sim/0,prod/3.25}
    \draw[evidenceline] (\repo.south) -- (\xpos,1.04);
  \draw[evidenceline] (0,1.04) -- (ci.north);
  \draw[flowarrow] (ci) -- (release);
\end{tikzpicture}
}
\end{center}

\column{0.35\textwidth}
\fixedpanel{5.30cm}{Controls to agree with DAQ Git admins}{
\begin{itemize}\tightitem
  \item Named maintainers, CODEOWNERS, interface owners.
  \item Protected branches and release rules.
  \item Required peer and interface-owner reviews.
  \item Versioned schemas and compatibility policy.
  \item Reuse DUNE DAQ release and CI conventions.
  \item Retained CI evidence and risk/decision records.
\end{itemize}
}
\end{columns}
\end{frame}

\begin{frame}[t]{Parallel lanes, explicit convergence gates}
\eyebrow{INDEPENDENT IMPLEMENTATION WHERE POSSIBLE; SHARED DECISIONS WHERE REQUIRED}
\vspace{1.2mm}

\begin{columns}[T,onlytextwidth]
\column{0.59\textwidth}
\begin{center}
\resizebox{0.98\linewidth}{!}{%
\begin{tikzpicture}[x=1cm,y=1cm]
  \node[
    systemcyan,
    text width=9.1cm,
    minimum height=0.85cm
  ] (baseline) at (4.7,2.45) {
    Shared baseline (Gate 0)\\
    governance + matrix v0 + schema registry + compatibility policy
  };

  \node[
    systempurple,
    text width=2.7cm,
    minimum height=2.05cm
  ] (boardslane) at (1.5,0.65) {
    \textbf{Lane A: board services}\\[1mm]
    DAPHNE hooks + actor\\
    WIB hooks + actor
  };
  \node[
    systemmint,
    text width=2.7cm,
    minimum height=2.05cm
  ] (clientlane) at (4.7,0.65) {
    \textbf{Lane B: clients + protection}\\[1mm]
    OPC-UA + DAQ clients\\
    rack power + DPS permit
  };
  \node[
    systemcyan,
    text width=2.7cm,
    minimum height=2.05cm
  ] (testlane) at (7.9,0.65) {
    \textbf{Lane C: test + deployment}\\[1mm]
    simulator + evidence\\
    infrastructure + deployment
  };

  \foreach \lane in {boardslane,clientlane,testlane}
    \draw[evidenceline]
      (baseline.south -| \lane.north) -- (\lane.north);

  \node[
    systembox,
    draw=warmyellow,
    fill=warmyellow!14!noirpanel,
    text width=9.1cm,
    minimum height=0.95cm
  ] (gate) at (4.7,-1.55) {
    Pilot convergence: authority + states + protection timing + recovery\\
    release BOM + crossing-fault evidence
  };
  \foreach \lane in {boardslane,clientlane,testlane}
    \draw[interfaceline]
      (\lane.south) -- (gate.north -| \lane.south);
\end{tikzpicture}
}
\end{center}

\column{0.38\textwidth}
\fixedpanel{5.2cm}{May proceed in parallel}{
\begin{itemize}\tightitem
  \item Lane A: DAPHNE and WIB hooks and board actors.
  \item Lane B: bridge/DAQ clients and rack-power/DPS path.
  \item Lane C: simulator, CI evidence, and deployment skeleton.
\end{itemize}
\vspace{1mm}
\textcolor{riskred}{\bfseries Must converge before writes/release}\\
Canonical schema ownership; command authority; allowed states; timing and
failure survival; independent readback; recovery; compatible artifact set.
}
\end{columns}
\end{frame}

\begin{frame}[t]{Begin with two vertical slices}
\eyebrow{SMALL ENOUGH TO CLOSE; BROAD ENOUGH TO EXPOSE SYSTEM DEPENDENCIES}
\vspace{1.0mm}

\begin{columns}[T,onlytextwidth]
\column{0.48\textwidth}
\fixedpanel{3.65cm}{1. Shared DAPHNE/WIB OPC-UA path}{
\begin{itemize}\tightitem
  \item Preserve native ZMQ/protobuf device contracts.
  \item Publish both board classes through one fleet bridge implementation.
  \item Exercise monitoring, guarded configuration, queueing, and readback.
  \item Trace \texttt{drunc}/appfwk init and conf dependencies.
  \item Measure bridge count, leaf concurrency, latency, error, and recovery.
\end{itemize}
}

\column{0.48\textwidth}
\fixedpanel{3.65cm}{2. Rack power and DPS-permit path}{
\begin{itemize}\tightitem
  \item PDU, WIENER, sensors, PLC/DPS permits, alarms, and latches.
  \item Normal energization and controlled shutdown sequence.
  \item Network loss, stale data, trip, reset, and recovery.
  \item Demonstrate that Ignition cannot bypass the independent permit.
  \item Capture owner-approved timing and failure evidence.
\end{itemize}
}
\end{columns}

\vspace{1mm}
\begin{beamercolorbox}[wd=\linewidth,sep=6pt,rounded=true]{panel alt}
\centering
\textcolor{purpletest}{\bfseries Crossing test:}\quad
inject one power, temperature, or communications fault that affects both
board configuration and rack protection, then verify the complete recovery path.
\end{beamercolorbox}
\end{frame}

\begin{frame}[t]{A practical Vertical Slice Test platform}
\eyebrow{DEVELOPMENT HARDWARE NOW; PRODUCTION SIZING FROM MEASUREMENT}
\vspace{1mm}

\begin{columns}[T,onlytextwidth]
\column{0.61\textwidth}
\begin{center}
\resizebox{0.93\linewidth}{!}{%
\begin{tikzpicture}[node distance=0.6cm and 0.62cm]
  \node[systemcyan,text width=3.25cm,minimum height=1.2cm] (host1) {
    Dell Precision 7680 \#1\\
    Linux Ignition Gateway\\32 GB RAM
  };
  \node[systempurple,text width=3.25cm,minimum height=1.2cm,right=of host1] (host2) {
    Dell Precision 7680 \#2\\
    bridges, load + faults\\32 GB RAM
  };
  \node[systemorange,text width=3.45cm,minimum height=1.2cm,below=0.9cm of host1,xshift=1.95cm] (switch) {
    Managed development switch\\
    VLANs, counters, mirror port
  };
  \node[systemmint,text width=2.55cm,below left=0.8cm and -0.15cm of switch] (plc) {
    PLC / DPS\\representative I/O
  };
  \node[systemmint,text width=2.55cm,right=0.45cm of plc] (power) {
    PDU + WIENER\\or safe emulators
  };
  \node[systemcyan,text width=2.55cm,right=0.45cm of power] (boards) {
    DAPHNE + WIB\\hardware or simulators
  };
  \draw[interfaceline] (host1) -- (switch);
  \draw[interfaceline] (host2) -- (switch);
  \draw[interfaceline] (switch) -- (plc);
  \draw[interfaceline] (switch) -- (power);
  \draw[interfaceline] (switch) -- (boards);
\end{tikzpicture}
}
\end{center}

\vspace{0.4mm}
\tinycite{The workstations are a development/VST resource, not a proposed SURF production-server specification.}

\column{0.36\textwidth}
\fixedpanel{5.15cm}{Inputs still needed}{
\begin{itemize}\tightitem
  \item SURF Ignition server, redundancy, storage, and virtualization specs.
  \item Controls network and managed-switch zones, routes, timing, and endpoint count.
  \item PLC/DPS platform, representative endpoints, and safe test-I/O boundary.
  \item Approved power/UPS/recovery cases and quantitative sizing criteria.
\end{itemize}
\vspace{0.6mm}
\textcolor{noircyan}{\bfseries Architecture rule}\\
Gateway, bridge-process, host, and failover counts are separate sizing decisions.
}
\end{columns}
\end{frame}

\begin{frame}[t]{Move quickly in the right order}
\eyebrow{A DISCUSSION STARTING POINT---DATES FOLLOW OWNERSHIP AND STAFFING}
\vspace{1.5mm}

\begin{tikzpicture}[x=1cm,y=1cm]
  \node[decisionbox] (now) at (1.25,0) {
    \textcolor{warmyellow}{\bfseries 1\enspace Now}\\[1mm]
    Agree governance, matrix authority, owners, and pilot boundaries.
  };
  \node[decisionbox] (v0) at (4.15,0) {
    \textcolor{warmyellow}{\bfseries 2\enspace Matrix v0}\\[1mm]
    Issue one template; record known contracts, conflicts, and missing inputs.
  };
  \node[decisionbox] (pilots) at (7.05,0) {
    \textcolor{warmyellow}{\bfseries 3\enspace Pilot closure}\\[1mm]
    Link packages, run VST cases, close two vertical paths with evidence.
  };
  \node[decisionbox] (schedule) at (9.95,0) {
    \textcolor{warmyellow}{\bfseries 4\enspace Staffed schedule}\\[1mm]
    Convert matrix dependencies into estimates, milestones, and release gates.
  };
  \node[decisionbox] (scale) at (12.7,0) {
    \textcolor{warmyellow}{\bfseries 5\enspace Scale-out}\\[1mm]
    Repeat network/server curves; extend only through reviewed rows.
  };
  \draw[draw=noircyan,line width=1.1pt] (0.25,-1.30) -- (13.48,-1.30);
  \fill[noircyan]
    (13.92,-1.30) -- (13.48,-1.05) -- (13.48,-1.55) -- cycle;
  \node[labelnote,above] at (7.1,-1.30)
    {dependency order and increasing implementation scope};
\end{tikzpicture}

\vspace{4mm}
\begin{columns}[T,onlytextwidth]
\column{0.48\textwidth}
\panelbox{Why this is faster}{
Unknowns are assigned early; integration work is ordered by dependency;
package changes receive immediate compatibility evidence.}
\column{0.48\textwidth}
\panelbox{Why this is safer}{
Unsupported guarantees and software-independent protection paths become
visible before commissioning, not during recovery from a failed test.}
\end{columns}

\vspace{2mm}
\centering
\color{noirmuted}\small
The first deliverable is not more disconnected code: it is matrix v0 with owners, sources, and open questions.
\end{frame}

\begin{frame}[t]{Decisions requested from this group}
\eyebrow{AUTHORIZE THE SHARED VERSIONED WORKSPACE, THEN THE WORK}
\vspace{1.0mm}

\begin{columns}[T,onlytextwidth]
\column{0.48\textwidth}
\fixedpanel{3.2cm}{Authorize Gate 0}{
\begin{enumerate}\tightitem
  \item \textbf{Create a collaboration-controlled SC/DPS Git group} with DAQ
    Git-admin support and named maintainers.
  \item Use the interaction matrix to set technical priorities and schedule.
  \item Name subsystem and lifecycle interface owners.
  \item Approve two initial vertical slices and a crossing-fault case.
\end{enumerate}
}

\column{0.48\textwidth}
\fixedpanel{3.2cm}{Provide the missing inputs}{
\begin{enumerate}\tightitem
  \item Consortium variables, commands, constraints, interlocks, and recovery contracts.
  \item DAQ Git-admin help creating the group and defining protections,
    releases, and retained CI evidence.
  \item SURF VST and production Ignition/server/network specifications.
  \item PLC, switch, PDU, power, UPS, and VST endpoint availability.
\end{enumerate}
}
\end{columns}

\vspace{1mm}
\begin{beamercolorbox}[wd=\linewidth,sep=6pt,rounded=true]{panel decision}
\centering
\color{noirbg}\bfseries\fontsize{11}{13}\selectfont
Gate 0: create the shared Git group and publish matrix v0.\\
That turns the simulator from isolated work into a versioned integration and
verification framework---and makes labor and completion estimates credible.
\end{beamercolorbox}
\end{frame}

\appendix

\begin{frame}[t]{Appendix: current quantitative evidence}
\eyebrow{LOCAL TEST BENCH AND SOURCE-DERIVED PLANNING VALUES}
\vspace{1.5mm}

\scriptsize
\begin{tabularx}{\linewidth}{@{}>{\bfseries}p{3.1cm}p{2.45cm}p{3.1cm}X@{}}
\toprule
Result & Value & Evidence class & Boundary \\
\midrule
Logical fleet & 1,818 objects & implemented + synthetic &
190 DAPHNE, 480 WIB, 108 PDU, 208 power channels, sensors, permits, alarms.\\
OPC-UA namespace & 40,755 variables & workstation test &
One shared fleet endpoint; local Linux host.\\
Balanced campaign & 12,000 attempts, 0 errors & workstation test &
3,000 trials/class; simulator time scale zero.\\
Bulk OPC-UA read & 0.516 s & workstation test &
Software snapshot, not a real detector scan guarantee.\\
OPC-UA process knee & 4 processes & workstation test &
Same physical host; read and address-space churn workload.\\
DAPHNE floor & 2.6901 s/board & source-derived projection &
Explicit waits; some hardware programming time remains unknown.\\
WIB floor & 16.01--16.10 s/board & source-derived projection &
COLDADC calibration/SPI waits; other hardware I/O remains unknown.\\
WIB fleet candidate & $4\times48$ leaves, 48.30 s floor &
posterior-predictive projection &
Four bridge queues, three waves; requires async configuration jobs.\\
\bottomrule
\end{tabularx}

\vspace{1.5mm}
\tinycite{Primary source: Far-Detector Rack-Scale Slow-Control Emulation and Interface Verification, DCE-FD2VD-2026-07-27.}
\end{frame}

\begin{frame}[t]{Appendix: terms that must not be conflated}
\eyebrow{ONE WORD---``SERVER''---HAD BEEN HIDING FOUR DIFFERENT DECISIONS}
\vspace{2mm}

\begin{columns}[T,onlytextwidth]
\column{0.24\textwidth}
\fixedpanel{3.35cm}{Ignition Gateway}{
The SCADA application layer: tags, alarms, history, sessions, and operator
views. One Gateway is the current development assumption.
}
\column{0.24\textwidth}
\fixedpanel{3.35cm}{Bridge instance}{
One OPC-UA endpoint/process translating native device protocols into a fleet
namespace. Four is the same-host monitoring candidate.
}
\column{0.24\textwidth}
\fixedpanel{3.35cm}{Leaf worker}{
One active in-flight board configuration job. Forty-eight per bridge is the
initial WIB planning point; it is not yet implemented.
}
\column{0.24\textwidth}
\fixedpanel{3.35cm}{Physical host}{
A Linux machine or VM carrying one or more processes. Host and failover counts
require a multi-host, representative-network test.
}
\end{columns}

\vspace{1mm}
\begin{beamercolorbox}[wd=\linewidth,sep=6pt,rounded=true]{panel alt}
\centering
\bfseries
Four bridge processes on one workstation is neither four Ignition Gateways nor four physical servers.
\end{beamercolorbox}
\end{frame}

\begin{frame}[t]{Appendix: interaction-matrix starter row}
\eyebrow{THE FORMAT SHOULD MAKE AN UNKNOWN IMPOSSIBLE TO HIDE}
\vspace{1.3mm}

\scriptsize
\setlength{\tabcolsep}{3.5pt}
\begin{tabularx}{\linewidth}{@{}p{1.55cm}p{1.35cm}p{1.65cm}p{1.65cm}p{1.55cm}p{1.65cm}p{1.55cm}X@{}}
\toprule
\textbf{Interaction} &
\textbf{Stage / state} &
\textbf{Authority} &
\textbf{Precondition} &
\textbf{Timing} &
\shortstack[l]{\textbf{Failure}\\\textbf{/ fallback}} &
\textbf{Owner} &
\textbf{Evidence / status} \\
\midrule
Configure one WIB &
commis\-sioning; powered; timing ready &
DAQ request through guarded bridge &
DPS permit; identity/config hash; FEMB power &
16.10 s explicit-wait floor; timeout TBD &
reject if permit lost; retry class TBD; independent readback &
WIB / DAQ / SC / DPS owners to confirm &
Serializer mapped; simulator exercised; full OPC-UA job queue and real hardware timing unresolved \\
\addlinespace
Rack power loss &
all stages; abnormal power &
independent protection path required &
approved shutdown ordering and hold-up assumptions &
APA-bias/WIB ordering bound TBD &
software may disappear; protection must remain effective &
electrical, APA-bias, WIB, SC/DPS owners to confirm &
Conditional risk identified; physical path, timing, and proof test unresolved \\
\bottomrule
\end{tabularx}

\vspace{1mm}
\begin{beamercolorbox}[wd=\linewidth,sep=4pt,rounded=true]{panel alt}
\centering
Rows begin as questions. They become schedule inputs only when the owner,
source revision, acceptance criterion, and evidence route are explicit.
\end{beamercolorbox}
\end{frame}

\begin{frame}[t]{Appendix: compare like with like}
\eyebrow{QUASAR WAS THE BRIDGE; THE LARGE LOAD INCLUDED EMULATORS AND IGNITION TAGS}
\vspace{0.4mm}

\scriptsize
\setlength{\tabcolsep}{4pt}
\begin{tabularx}{\linewidth}{
  @{}>{\bfseries}p{2.20cm}p{5.55cm}X@{}}
\toprule
 & \textbf{March 2026 WIB/Quasar test} &
   \textbf{July 2026 local fleet bench} \\
\midrule
Runtime shape &
Many Python WIB emulator processes/endpoints
\(\rightarrow\) ZMQ \(\rightarrow\) Quasar &
1,818 in-process logical objects
\(\rightarrow\) shared async runtime.\\
Namespace / tags &
At 1,000 WIBs: \(\sim5.5\) million Ignition tags; only a small subset were
OPC tags. &
40,755 OPC-UA variables across the full fleet; compact read-only Ignition
resources.\\
\shortstack[l]{Observed\\resources} &
At 480 WIBs: emulator 21.7 GB; Quasar 3.23 GB; Ignition 55.6 GB and
\(\sim6.7\) CPU cores. &
Emulator 0.65 GiB; DAPHNE bridge 0.05 GiB; Ignition 1.6 GiB.\\
Test boundary &
Large connection and Ignition tag/alarm stress test. &
Loopback interface test; production scan classes, history, and client load not
reproduced.\\
\bottomrule
\end{tabularx}

\vspace{0mm}
\begin{beamercolorbox}[wd=\linewidth,sep=3pt,rounded=true]{panel alt}
{\bfseries Controlled comparison: hold the workload fixed and vary one layer\par}
\vspace{0.35mm}
\scriptsize
\begin{tabularx}{\dimexpr\linewidth-10pt\relax}{
  @{}>{\bfseries}p{2.15cm}X@{}}
Hold constant &
480 identities and protobuf; OPC-UA nodes; Ignition UDTs/tags/alarms; scan,
history, clients, and fault script.\\
Experiment A &
Same bridge and Ignition project; vary process-per-WIB versus shared-async
emulation.\\
Experiment B &
Same emulator endpoints and Ignition project; vary Quasar versus the local
bridge.\\
Measure &
Per-component CPU/RSS, descriptors, connections, cold start, throughput,
p95/p99 latency, tag age/quality, errors, and reconnect recovery.\\
\end{tabularx}
\vspace{0.2mm}
\centering\bfseries
Only then compare efficiency against agreed latency, recovery, and resource
headroom.
\end{beamercolorbox}

\vspace{0mm}
\tinycite{\centering Sources: DUNE Slow Controls ZMQ/OPC-UA Spin Test,
12 March 2026; DCE-FD2VD-2026-07-27.}
\end{frame}

\end{document}
